\documentclass[letterpaper,10pt]{article}

\usepackage{latexsym}
\usepackage[empty]{fullpage}
\usepackage{titlesec}
\usepackage{marvosym}
\usepackage[usenames,dvipsnames]{color}
\usepackage{verbatim}
\usepackage{enumitem}
\usepackage[hidelinks]{hyperref}
\usepackage{fancyhdr}
\usepackage[english]{babel}
\usepackage{tabularx}
\input{glyphtounicode}

\pagestyle{fancy}
\fancyhf{} % clear all header and footer fields
\fancyfoot{}
\renewcommand{\headrulewidth}{0pt}
\renewcommand{\footrulewidth}{0pt}

% Adjust margins
\addtolength{\oddsidemargin}{-0.5in}
\addtolength{\evensidemargin}{-0.5in}
\addtolength{\textwidth}{1in}
\addtolength{\topmargin}{-.5in}
\addtolength{\textheight}{1.0in}

\urlstyle{same}

\raggedbottom
\raggedright
\setlength{\tabcolsep}{0in}

% Sections formatting
\titleformat{\section}{
  \vspace{-4pt}\scshape\raggedright\large
}{}{0em}{}[\color{black}\titlerule \vspace{-5pt}]

% Ensure that generate pdf is machine readable/ATS parsable
\pdfgentounicode=1

%-------------------------
% Custom commands
\newcommand{\resumeItem}[1]{
  \item\small{
    {#1 \vspace{-2pt}}
  }
}

\newcommand{\resumeSubheading}[4]{
  \vspace{-2pt}\item
    \begin{tabular*}{0.97\textwidth}[t]{l@{\extracolsep{\fill}}r}
      \textbf{#1} & #2 \\
      \textit{\small#3} & \textit{\small #4} \\
    \end{tabular*}\vspace{-7pt}
}

\newcommand{\resumeSubSubheading}[2]{
    \item
    \begin{tabular*}{0.97\textwidth}{l@{\extracolsep{\fill}}r}
      \textit{\small#1} & \textit{\small #2} \\
    \end{tabular*}\vspace{-7pt}
}

\newcommand{\resumeProjectHeading}[2]{
    \item
    \begin{tabular*}{0.97\textwidth}{l@{\extracolsep{\fill}}r}
      \small#1 & #2 \\
    \end{tabular*}\vspace{-7pt}
}

% \newcommand{\resumeSubItem}[1]{\resumeItem{#1}\vspace{-4pt}}

\renewcommand\labelitemii{$\vcenter{\hbox{\tiny$\bullet$}}$}

\newcommand{\resumeSubHeadingListStart}{\begin{itemize}[leftmargin=0.15in, label={}]}
\newcommand{\resumeSubHeadingListEnd}{\end{itemize}}
\newcommand{\resumeItemListStart}{\begin{itemize}}
\newcommand{\resumeItemListEnd}{\end{itemize}\vspace{-5pt}}

%-------------------------------------------
%%%%%%  RESUME STARTS HERE  %%%%%%%%%%%%%%%%%%%%%%%%%%%%

% Add publications
% \usepackage{etoolbox}
\usepackage[
style=numeric,
sorting=ydnt,
defernumbers=true,
maxbibnames=99,
maxcitenames=99
]{biblatex}

% \AtDataInput{%
%   \csnumgdef{entrycount:\therefsection}{%
%     \csuse{entrycount:\therefsection}+1}}

% \DeclareFieldFormat{labelnumber}{\mkbibdesc{#1}}    
% \newrobustcmd*{\mkbibdesc}[1]{%
%   \number\numexpr\csuse{entrycount:\therefsection}+1-#1\relax}

\renewcommand*{\mkbibcompletename}[1]{%
  \ifitemannotation{highlight}
    {\textbf{#1}}
    {#1}%
  \ifitemannotation{first}
    {\kern-0.05em\textsuperscript{$^*$}}
    {}%
  \ifitemannotation{last}
    {\textsuperscript{$\ddagger$}}
    {}%
}

\addbibresource[label=preprints]{preprints.bib}
\addbibresource[label=accepted]{accepted.bib}
\addbibresource[label=publications]{publications.bib}

\begin{document}
\begin{center}
    \textbf{\Huge \scshape Jun Zeng} \\ \vspace{2pt}
    
    \small zengjunsjtu@berkeley.edu $|$ 
    \href{https://junzengx14.github.io/}{https://junzengx14.github.io/} $|$
    \href{https://scholar.google.com/citations?hl=en&user=xBOE358AAAAJ}{Google Scholar} $|$
    \href{https://github.com/junzengx14}{GitHub}
\end{center}

The latest update of CV was done by Jun Zeng on \today.

\section{Research Experience}
I am interested in researching mathematical foundations and developing modern algorithms for
    autonomous robots, achieving low-latency, agile, and safe maneuvers with adaptation and robustness
    in complex environments. My research lies in the combination of control, optimization, motion
    planning, trajectory generation, reinforcement learning, and machine learning.

\section{Experience}
    \resumeSubHeadingListStart
        \resumeSubheading
        {Chief Technology Officer}{Oct 2025 -- Present}
        {Cyan Robotics}{Shanghai, China}
        \resumeItemListStart
            \resumeItem{Technical Leadership: Own technology strategy, system architecture, and
                engineering roadmap. Grew engineering from $\sim$15 to $\sim$40 across mechanical and
                electrical design, embedded systems, robot control and whole-body autonomy, and
                multimodal end-to-end models, keeping hardware and software teams on a single roadmap.}
            \resumeItem{Dino OS: Proposed and architected a modern embodied AI framework unifying
                build, test, training, fine-tuning and deployment across heterogeneous compute platforms. 
                One reproducible toolchain from experiment to on-robot deployment, and the
                pivotal foundation for a fast-moving, complex embodied program; deployed on Amoo.}
            \resumeItem{Hardware-Software Co-Design: Directed in-house IMU, stereo dToF module, and
                modular robot joints. Vertical integration cut BOM cost and removed vendor lock-in,
                while native driver integration enabled cross-sensor time synchronization and low-level
                control unavailable off the shelf.}
            \resumeItem{Product \& Business Impact: Took Amoo from research prototype to integrated
                consumer product. Global debut at AWE 2026, followed by various venues such as BEYOND Expo 2026 in Macao,
                WAIC in Shanghai, WRC in Beijing, IFA in Berlin; over
                1000 letters of intent to purchase.}
        \resumeItemListEnd

        \resumeSubheading
        {Senior AI/ML Scientist}{Jun 2022 -- Oct 2025}
        {Cruise LLC}{San Francisco, CA}
        \resumeItemListStart
            \resumeItem{Large Behavior Models \& AI Guardrails: Worked as one of core ICs for end-to-end large
                behavior model experience spanning pretraining, post-training, evaluation harness, and
                deployment. Designed and improved auto-regressive trajectory decoder,
                enhancing dynamic feasibility of the generated trajectories from the large ML model.
                Designed and contributed improvements to trajectory constraints enforcement in both
                training and inference time using various techniques such as flow matching, adapter
                layers fine-tuning, etc. Contributed to improve algorithm for training stability and
                robustness when transitioning between manual and autonomous driving, ensuring better
                policy learning and deployment.}
            \resumeItem{Trajectory Generation: Led the unification of reactive motion planning and
                control to enhance real-time adaptability and reduce system complexity. Led next
                generation of coupled non-convex trajectory generation solver with latest techniques of
                inner point methods. Designed and developed ML-based warm start calculation to
                accelerate solver runtime. Designed and integrated a ML-based dragging detector with
                collision imminence.}
            \resumeItem{Cross Team Technical Leadership: Worked as one of core ICs to reduce
                unnecessary braking and swerving to improve company-level top-down metrics; Collaborated
                closely with prediction and behavior planning teams to refine trajectory generation for
                complex driving scenarios, such as tight turns and cautious VRU interactions,
                contributing to more robust large planner model training and ML-based trajectory
                selection.}
        \resumeItemListEnd
    \resumeSubHeadingListEnd

\section{Education}
\hspace*{0.5cm} \textbf{University of California, Berkeley} \hfill Aug 2017 - May 2022 \\
\hspace*{0.5cm} Ph.D. in Mechanical Engineering \\
\hspace*{0.5cm} Advisor: Prof. Koushil Sreenath \\
\hspace*{0.5cm} Field: Control, Robotics, Machine Learning
\break

\hspace*{0.5cm} \textbf{Ecole Polytechnique} \hfill Mar 2015 - Aug 2017 \\
\hspace*{0.5cm} Diplôme d'ingénieur in Applied Mathematics and Mechanics
\break

\hspace*{0.5cm} \textbf{Shanghai Jiao Tong University} \hfill Sep 2012 - May 2016 \\
\hspace*{0.5cm} Bachelor of Science and Engineering in Automation

\section{Honors \& Awards}
\begin{itemize}
    \item \textbf{Chin Leung Shui Chun Fellowship}, UC Berkeley \hfill \emph{2021}
    \item \vspace*{-0.25cm} \textbf{Graduate Division Block Fellowship}, UC Berkeley \hfill \emph{2020 - 2021}
    \item \vspace*{-0.25cm} \textbf{Mechanical Engineering Department Award}, UC Berkeley \hfill \emph{2020}
    \item \vspace*{-0.25cm} \textbf{X Fondation Fellowship}, Ecole Polytechnique \hfill \emph{2015 - 2017}
    \item \vspace*{-0.25cm} \textbf{Ardian Excellence Scholarship}, SJTU \hfill \emph{2014}
    \item \vspace*{-0.25cm} \textbf{University Excellence Scholarship}, SJTU \hfill \emph{2012 - 2014}
\end{itemize}

% \begin{refsection}[preprints]
% \nocite{*}
% \printbibliography[title = {Preprints}]
% \end{refsection}

% \begin{refsection}[accepted]
% \nocite{*}
% \printbibliography[title = {Accepted Work}]
% \end{refsection}

\begin{refsection}[publications]
\nocite{*}
\printbibliography[title = {Publications}]
\end{refsection}

\section{Open-Source Tools \& Software}
\hspace*{0.5cm}
\begin{tabularx}{0.95\textwidth}{X}
\textbf{Car-Racing: A toolkit to validate algorithm for Autonomous Car Racing} \\
- Versatile and Scalable simulation platform for autonomous car racing problems supporting dynamics,
    control, planning algorithm. \\
- Code repository: 
\href{https://github.com/HybridRobotics/car-racing}{\textit{https://github.com/HybridRobotics/car-racing}}
\end{tabularx}
\\ \hfill \break

\hspace*{0.5cm}
\begin{tabularx}{0.95\textwidth}{X}
\textbf{Python function for APIs for control barrier functions in the field of control, planning and navigation} \\
- A python API library of multiple algorithm of control barrier functions for control and robotics
    applications. \\
- Code repository: \href{https://github.com/HybridRobotics/cbf}{\textit{https://github.com/HybridRobotics/cbf}}
\end{tabularx}
\\ \hfill \break

\hspace*{0.5cm}
\begin{tabularx}{0.95\textwidth}{X}
\textbf{Code Collection of Discrete-time Control Barrier Functions (DCBF)} \\
- A collection of multiple algorithm of discrete-time control barrier functions for control and
    robotics applications. This collection makes it simple to track latest work about control barrier
    functions. \\
- Code repository: \href{https://github.com/HybridRobotics/NMPC-DCLF-DCBF}{\textit{https://github.com/HybridRobotics/NMPC-DCLF-DCBF}}
\end{tabularx}

\section{Invited Talks}
\hspace*{0.5cm} \textbf{Safety-Critical Model Predictive Control with Discrete-Time Control Barrier Function} \\
\hspace*{0.5cm} \emph{3rd Norcal Control Workshop, Stanford University (Virtual), CA} \hfill Jan 2021

\section{Professional Service}

\hspace*{0.5cm} \textbf{Program Committee} \\
\hspace*{0.5cm} - AAAI-23 special track on Safe and Robust AI \hfill \emph{2023} \\
\hspace*{0.5cm} - ICRA-2022 Workshop on Opportunities and Challenges with Autonomous Racing (OCAR-ICRA)  \hfill \emph{2022} \\

\hspace*{0.5cm} \textbf{Journal Reviewers} \\
\hspace*{0.5cm} - ASME Journal of Dynamic Systems, Measurement, and Control (J. Dyn. Sys., Meas., Control.) \hfill \emph{2021 - 2025} \\
\hspace*{0.5cm} - Automatica \hfill \emph{2022 - 2026} \\
\hspace*{0.5cm} - Control Engineering Practice \hfill \emph{2025 - 2026} \\
\hspace*{0.5cm} - IEEE Access \hfill \emph{2022 - 2025} \\
\hspace*{0.5cm} - IEEE Control Systems Letters (L-CSS) \hfill \emph{2023 - 2026} \\
\hspace*{0.5cm} - IEEE Open Journal of Control Systems (OJCSYS) \hfill \emph{2023 - 2025} \\
\hspace*{0.5cm} - IEEE Robotics and Automation Letters (RA-L) \hfill \emph{2020 - 2026} \\
\hspace*{0.5cm} - IEEE Transactions on Automatic Control (TAC) \hfill \emph{2021 - 2026} \\
\hspace*{0.5cm} - IEEE Transactions on Automation Science and Engineering (T-ASE) \hfill \emph{2026} \\
\hspace*{0.5cm} - IEEE Transactions on Control of Network Systems (T-CONES) \hfill \emph{2021 - 2025} \\
\hspace*{0.5cm} - IEEE Transactions on Control Systems Technology (T-CST) \hfill \emph{2021 - 2026} \\
\hspace*{0.5cm} - IEEE Transactions on Industrial Electronics (TIE) \hfill \emph{2024 - 2026} \\
\hspace*{0.5cm} - IEEE Transactions on Intelligent Transportation Systems (T-ITS) \hfill \emph{2024 - 2026} \\
\hspace*{0.5cm} - IEEE Transactions on Robotics (TR-O) \hfill \emph{2023 - 2026} \\
\hspace*{0.5cm} - IEEE Transactions on Systems, Man, and Cybernetics (TSMC) \hfill \emph{2026} \\
\hspace*{0.5cm} - IEEE/ASME Transactions on Mechatronics (T-Mech) \hfill \emph{2022 - 2026} \\
\hspace*{0.5cm} - International Journal of Robotics Research (IJRR) \hfill \emph{2023 - 2025} \\
\hspace*{0.5cm} - International Journal of Robust and Nonlinear Control (RNC) \hfill \emph{2023 - 2025} \\

\hspace*{0.5cm} \textbf{Conference Reviewers} \\
\hspace*{0.5cm} - AAAI Conference on Artificial Intelligence (AAAI) \hfill \emph{2022 - 2024} \\
\hspace*{0.5cm} - American Control Conference (ACC) \hfill \emph{2023 - 2026} \\
\hspace*{0.5cm} - Annual Conference of the IEEE Industrial Electronics Society (IECON) \hfill \emph{2023 - 2025} \\
\hspace*{0.5cm} - European Control Conference (ECC) \hfill \emph{2023 - 2025} \\
\hspace*{0.5cm} - IEEE Conference on Control Technology and Applications (CCTA) \hfill \emph{2024 - 2025} \\
\hspace*{0.5cm} - IEEE Conference on Decision and Control (CDC) \hfill \emph{2022 - 2026} \\
\hspace*{0.5cm} - IEEE International Conference on Robotics and Automation (ICRA) \hfill \emph{2020 - 2026} \\
\hspace*{0.5cm} - IEEE International Conference on Unmanned Aircraft Systems (ICUAS) \hfill \emph{2020 - 2021} \\
\hspace*{0.5cm} - IEEE/RSJ International Conference on Intelligent Robots and Systems (IROS) \hfill \emph{2021 - 2025} \\
\hspace*{0.5cm} - Learning for Dynamics \& Control Conference (L4DC) \hfill \emph{2023 - 2026} \\

\section{Other Professional Experiences}
\hspace*{0.5cm} \textbf{Waymo LLC} \hfill May 2021 - Aug 2021 \\
\hspace*{0.5cm} Software Engineer Intern \\
\hfill

\hspace*{0.5cm} \textbf{PlusAI} \hfill May 2019 - Aug 2019 \\
\hspace*{0.5cm} Software Engineer Intern \\
\hfill

\hspace*{0.5cm} \textbf{Valeo} \hfill May 2016 - Aug 2016 \\
\hspace*{0.5cm} Summer Intern \\
\hfill

\hspace*{0.5cm} \textbf{Autonomous Systems and Robotics Lab, ENSTA} \hfill Sep 2015 -- Mar 2016 \\
\hspace*{0.5cm} Research Assistant \\

\section{Teaching Experience}
\begin{itemize}
    \item \textbf{EECS 106B/206B:  Robotic Manipulation and Interaction} \hfill 2020 Spring \\ Graduate Student Instructor, UC Berkeley \\
    \item \textbf{EECS 106A/206A:  Introduction to Robotics} \hfill 2019 Fall \\ Graduate Student Instructor, UC Berkeley
\end{itemize}

% \section{Research Mentoring}
% \begin{itemize}
%     \item Lizhi Yang \hfill \emph{2019 - 2022} \\
%     \emph{undergraduate at UC Berkeley, now Ph.D. student at California Institute of Technology} \\
%     Worked on legged robotics, work published at ICRA 2021/2022, CASE 2021 and RA-L 2022
%     \item Suiyi He \hfill \emph{2020 - 2021} \\
%     \emph{MEng Student at UC Berkeley, now Ph.D. student at University of Minnesota Twin Cities} \\
%     Worked on autonomous driving and racing, work published at ACC 2021 and ICRA 2022
%     \item Anxing Xiao \hfill \emph{2020 - 2021} \\
%     \emph{visiting undergraduate, now Ph.D. Student at University of British Columbia} \\
%     Worked on manipulation problems with legged robots, work published at ICRA 2021
%     \item Chenyu Yang \hfill \emph{2020 - 2022} \\
%     \emph{visiting undergraduate, now MS Student at ETH Zurich} \\
%     Work on legged robotics manipulation problems, work published at IROS 2020 and RA-L 2022
% \end{itemize}

% \section{Language \& Skills}
% \begin{itemize}
%     \item \textbf{Programming}: Linux Shell, C/C++, Python, SQL, JavaScript, TypeScript.
%     \item \textbf{Tools}: Bazel, CMake.
%     \item \textbf{Robotics}: ROS.
%     \item \textbf{Optimization}: OSQP, Ipopt, CVX, Jax, CasADi, Gurobi.
%     \item \textbf{Learning}: PyTorch, TensorFlow (Python/C++), Flax.
%     \item \textbf{Vision}: OpenCV, PCL, OpenGL.
% \end{itemize}

\end{document}
